\chapter{A Legal Example: Lawful, Right, Known, Required}
\label{ch:legal-example}

The Kerberos discussion needed two questions beyond temporal logic: \emph{who
knows which values}, and \emph{who is permitted to hold them}. Law adds more
pressure on the same machinery. Obligations have deadlines; they are triggered
by what someone knew; breaking one creates another; and two codes can disagree.
This chapter works one small legal example end to end so that each modality is
used for something concrete. Its tables and diagram are generated by
\path{LegalDepositExample.lean} from the event histories, not drawn by hand.

The statute below is invented, and so is the moral code. Neither describes any
real jurisdiction, and nothing here is legal advice.

\section{Three meanings of ``legal''}

LeanFM already uses the word \emph{legal} for histories that the grammar
admits. A legal example needs two further meanings, and conflating them is the
most common error in normative reasoning about software.

\begin{center}
\begin{tabular}{@{}p{0.17\textwidth}p{0.36\textwidth}p{0.37\textwidth}@{}}
\toprule
Term & Judged against & Operators \\
\midrule
admissible & the protocol grammar: what can happen at all & CTL $A$, $E$, $F$, $G$ over generated futures \\
lawful / unlawful & a declared statute's ideal futures & $O^{Act}_a$, $P^{Act}_a$ \\
right / wrong & a declared moral code with a named author & $O^{Code}_a$, $P^{Code}_a$ \\
\bottomrule
\end{tabular}
\end{center}

In this chapter ``admissible'' replaces ``legal'' wherever the grammar is meant.
An unlawful act is usually admissible: people can break the law, so
$EF\,\mathit{unlawful}$ is normally true, and that is not a defect of the model.
Statutes constrain which futures are \emph{ideal}, not which futures are
\emph{possible}. A moral code is handled with exactly the same machinery and a
different declared relation. LeanFM never decides what is right; it checks
histories against a code that someone wrote down and signed, and it shows where
codes disagree.

\section{The Model Tenancy Deposit Act and a declared code}

The actors are a tenant $T$, a landlord $L$, a deposit registry, the post and a
court. Writing $\mathit{Ideal}^{N}_a(h)$ for the futures that code $N$ regards as
ideal for actor $a$ after history $h$, the obligation and permission operators of
the previous chapter are indexed by the code:
\AccessibleEquation{Under code N actor a is obligated to ensure psi at history h if and only if psi holds on every future that code N regards as ideal for a at h. a is permitted to ensure psi if psi holds on some such future.}{\[
O^{N}_a\psi(h)\iff\forall\pi\in\mathit{Ideal}^{N}_a(h).\;\psi(\pi),
\qquad
P^{N}_a\psi(h)\iff\exists\pi\in\mathit{Ideal}^{N}_a(h).\;\psi(\pi).
\]}
Read: ``Under code $N$, $a$ must make $\psi$ true when every future that $N$
calls ideal does so, and may make $\psi$ true when at least one does.''
Prohibition is obligation of the negation: $\mathit{Forbidden}^{N}_a\psi \equiv
O^{N}_a\neg\psi$. $F_{\le d}$ is a bounded path operator: within $d$ days.

\begin{description}[leftmargin=1.6em]
\item[Act \S1 (protection).] Within 30 days of receiving a deposit, the landlord
registers it.
\AccessibleEquation{Always. if a deposit has been paid then under the Act the landlord is obligated within 30 days to register it.}{\[
AG\,(\mathit{paid}\rightarrow O^{Act}_L\,F_{\le 30}\,\mathit{protected})
\]}
\item[Act \S2 (return, triggered by knowledge).] Once the lease has ended
\emph{and the landlord knows the tenant's forwarding address}, the landlord
returns the deposit or itemises deductions within 14 days.
\AccessibleEquation{Always. if the lease has ended and the landlord knows the forwarding address then under the Act the landlord is obligated within 14 days to return the deposit or itemise deductions.}{\[
AG\,\bigl((\mathit{ended}\wedge \mathit{knows}(L,\mathit{address}))\rightarrow
O^{Act}_L\,F_{\le 14}\,(\mathit{returned}\vee\mathit{itemised})\bigr)
\]}
\item[Act \S3 (reparation).] If \S2 is violated, the landlord pays the deposit
plus an equal penalty.
\AccessibleEquation{Always. if section 2 has been violated then under the Act the landlord is obligated eventually to pay the penalty.}{\[
AG\,(\mathit{violated}_{\S2}\rightarrow O^{Act}_L\,F\,\mathit{penaltyPaid})
\]}
\item[Act \S4 (prohibition).] A deposit is never applied to rent.
\AccessibleEquation{Under the Act it is obligatory that the deposit is always not applied to rent. equivalently applying it is forbidden.}{\[
O^{Act}\,G\,\neg\mathit{offset}\;\equiv\;\mathit{Forbidden}^{Act}\,F\,\mathit{offset}
\]}
\item[Code (declared by the stakeholder).] (a) Never charge for normal wear and
tear. (b) A landlord who holds another way to reach the tenant makes a
reasonable effort to use it. (c) Help a tenant in documented hardship when it
costs nothing you are owed.
\end{description}

\section{Knowledge from the causal past}

Act \S2 is conditioned on what the landlord knew. In LeanFM that is computed,
not asserted. Every event carries \texttt{prevIds} and \texttt{causeIds}; an
actor knows a value at an event when a message carrying it was delivered to
that actor somewhere in the event's causal ancestry:
\AccessibleEquation{Actor a knows value v at event e if and only if some event e prime in the causal past of e or e itself delivered to a a message from which v can be derived.}{\[
\mathit{knows}(a,v)@e \iff \exists e'\in \mathit{past}(e)\cup\{e\}.\;
\mathit{dst}(e')=a \wedge v\in\mathit{derive}(\mathit{msg}(e'))
\]}
Read: ``$a$ knows $v$ at event $e$ exactly when some event in $e$'s causal past
delivered to $a$ a message from which $v$ can be derived.'' This is
\emph{possession} knowledge, the kind security protocols need. The epistemic
operator $K_a$ of the previous chapter, which quantifies over histories that $a$
cannot tell apart, is reserved for claims such as ``the landlord could not have
known'' that range over alternatives.

The \emph{lost in the post} history shows why the distinction matters. The
tenant posts the address; the post never delivers it. The tenant then records
that she \emph{believes} \S2 was violated. She is wrong: the landlord's causal
past never contains the address, so \S2 is never triggered. Belief is an event
in the tenant's own history and is not factive; knowledge is computed and is.
The same history is nonetheless \emph{wrong} under Code (b): the landlord held
the tenant's email from the lease (event e1 is in every causal past) and made no
effort to use it.

\begin{figure}[ht]
\centering
\AccessibleFigure{Event table for the lost-in-the-post history. lease deposit protection and lease end then the tenant posts the address to the post office which never delivers it and finally the tenant records that she believes the Act was violated.}{\resizebox{\textwidth}{!}{% Generated by LegalDepositExample.lean; do not edit.
\begin{tabular}{@{}rrllp{0.3\textwidth}ll@{}}\toprule
Id & Day & From $\to$ To & Message & Detail & prevIds & causeIds\\\midrule
e1 & 0 & Landlord $\to$ Tenant & \texttt{LeaseSigned} & rent 1000, deposit 1000, tenant email on file &  & \\
e2 & 1 & Tenant $\to$ Landlord & \texttt{DepositPaid} & 1000 & e1 & e1\\
e3 & 12 & Landlord $\to$ Registry & \texttt{DepositProtected} & scheme ref 7731 & e2 & e2\\
e4 & 365 & Tenant $\to$ Landlord & \texttt{LeaseEnded} & keys returned & e2 & \\
e5 & 366 & Tenant $\to$ Post & \texttt{AddressPosted} & letter to landlord & e4 & \\
e6 & 400 & Tenant $\to$ Tenant & \texttt{BelievesViolated} & no refund after 34 days & e5 & \\
\bottomrule\end{tabular}
}}
\caption{The \emph{lost in the post} history. No event delivers \texttt{ForwardingAddress} to the landlord, so $\mathit{knows}(L,\mathit{address})$ never holds and Act~\S2 is never triggered.}
\end{figure}

\section{Violation and reparation}

A naive deontic logic breaks on the late return. If the landlord ought to
return the deposit by day~380, and ought to pay a penalty if he does not, then
after day~380 he seems to be both obliged to have returned it and obliged to pay
for not having done so. This is the contrary-to-duty problem, best known as
Chisholm's paradox.

Ideal futures dissolve it, because they are computed \emph{from the current
history}. Before day~380 every ideal future contains a timely return. After the
deadline passes, no admissible future contains one, so the ideal futures from
that history are the best ones still admissible: those in which the penalty is
paid. The obligation set changes as the history grows; it never becomes
inconsistent. The checker therefore reports three outcomes, not two:
\emph{fulfilled}, \emph{violated}, and \emph{violated, then repaired}.

\begin{figure}[ht]
\centering
\AccessibleFigure{Interaction diagram generated from the late-return history. Messages in order. day 0 lease signed from landlord to tenant. day 1 deposit paid from tenant to landlord. day 12 deposit protected from landlord to registry. day 365 lease ended from tenant to landlord. day 366 forwarding address from tenant to landlord. day 391 deposit returned from landlord to tenant. day 395 claim filed from tenant to court. day 410 judgment from court to landlord. day 420 penalty paid from landlord to tenant.}{\resizebox{\textwidth}{!}{% Generated by LeanFM.MDPDiagram.renderInteractionTeX.
\begin{tikzpicture}[actor/.style={draw=diagramink,text=diagramink,fill=diagramsurface,rounded corners,minimum width=2.3cm,minimum height=7mm,font=\small},message/.style={->,>=Stealth,draw=leanblue,thick}]
\node[actor] (a0) at (0,0) {Tenant};
\draw[dashed,draw=diagramink] (a0.south) -- ++(0,-10);
\node[actor] (a1) at (6,0) {Landlord};
\draw[dashed,draw=diagramink] (a1.south) -- ++(0,-10);
\node[actor] (a2) at (12,0) {Registry};
\draw[dashed,draw=diagramink] (a2.south) -- ++(0,-10);
\node[actor] (a3) at (18,0) {Court};
\draw[dashed,draw=diagramink] (a3.south) -- ++(0,-10);
\draw[message] (6,-1) -- node[above,fill=pagebackground,text=diagramink,font=\scriptsize] {e1 day 0: LeaseSigned} (0,-1);
\draw[message] (0,-2) -- node[above,fill=pagebackground,text=diagramink,font=\scriptsize] {e2 day 1: DepositPaid} (6,-2);
\draw[message] (6,-3) -- node[above,fill=pagebackground,text=diagramink,font=\scriptsize] {e3 day 12: DepositProtected} (12,-3);
\draw[message] (0,-4) -- node[above,fill=pagebackground,text=diagramink,font=\scriptsize] {e4 day 365: LeaseEnded} (6,-4);
\draw[message] (0,-5) -- node[above,fill=pagebackground,text=diagramink,font=\scriptsize] {e5 day 366: ForwardingAddress} (6,-5);
\draw[message] (6,-6) -- node[above,fill=pagebackground,text=diagramink,font=\scriptsize] {e6 day 391: DepositReturned} (0,-6);
\draw[message] (0,-7) -- node[above,fill=pagebackground,text=diagramink,font=\scriptsize] {e7 day 395: ClaimFiled} (18,-7);
\draw[message] (18,-8) -- node[above,fill=pagebackground,text=diagramink,font=\scriptsize] {e8 day 410: Judgment} (6,-8);
\draw[message] (6,-9) -- node[above,fill=pagebackground,text=diagramink,font=\scriptsize] {e9 day 420: PenaltyPaid} (0,-9);
\end{tikzpicture}
}}
\caption{The \emph{late return} history, rendered by \texttt{MDPDiagram.renderInteractionTeX} from its events. The address reaches the landlord on day~366, so the \S2 deadline is day~380; the return on day~391 violates it, and the penalty on day~420 repairs it.}
\end{figure}

\begin{figure}[ht]
\centering
\AccessibleFigure{Event table for the late-return history with each event's day sender receiver message detail prevIds and causeIds.}{\resizebox{\textwidth}{!}{% Generated by LegalDepositExample.lean; do not edit.
\begin{tabular}{@{}rrllp{0.3\textwidth}ll@{}}\toprule
Id & Day & From $\to$ To & Message & Detail & prevIds & causeIds\\\midrule
e1 & 0 & Landlord $\to$ Tenant & \texttt{LeaseSigned} & rent 1000, deposit 1000, tenant email on file &  & \\
e2 & 1 & Tenant $\to$ Landlord & \texttt{DepositPaid} & 1000 & e1 & e1\\
e3 & 12 & Landlord $\to$ Registry & \texttt{DepositProtected} & scheme ref 7731 & e2 & e2\\
e4 & 365 & Tenant $\to$ Landlord & \texttt{LeaseEnded} & keys returned & e2 & \\
e5 & 366 & Tenant $\to$ Landlord & \texttt{ForwardingAddress} & 12 Elm St & e4 & \\
e6 & 391 & Landlord $\to$ Tenant & \texttt{DepositReturned} & 1000 & e5, e3 & e5\\
e7 & 395 & Tenant $\to$ Court & \texttt{ClaimFiled} & Act s2 deadline day 380 missed & e6 & e6\\
e8 & 410 & Court $\to$ Landlord & \texttt{Judgment} & pay penalty 1000 under s3 & e7 & e7\\
e9 & 420 & Landlord $\to$ Tenant & \texttt{PenaltyPaid} & 1000 & e8 & e8\\
\bottomrule\end{tabular}
}}
\caption{The same history as events. \texttt{causeIds} pair each response with its cause (the return answers the address; the judgment answers the claim); \texttt{prevIds} order each actor's own events.}
\end{figure}

\section{Lawful, right, and the four quadrants}

\Cref{tab:legal-status} evaluates every rule in every history.
\Cref{tab:legal-verdicts} summarises each history under each code. All four
combinations of lawful and right appear, which is the point of keeping the codes
separate:

\begin{itemize}
\item \emph{lawful but wrong}: an on-time return that deducts for normal wear
(Code a), and the lost address the landlord could have chased (Code b);
\item \emph{violated, then repaired}: the late return, under the Act, with no
moral duty arising;
\item \emph{unlawful but right}: the landlord who, on a laid-off tenant's request,
applied the deposit to the final month's rent. That breaks Act \S4 and fulfils
Code (c).
\end{itemize}

\begin{table}[ht]
\centering
\small
\resizebox{\textwidth}{!}{% Generated by LegalDepositExample.lean; do not edit.
\begin{tabular}{@{}llllllll@{}}\toprule
History & Act \S1 & Act \S2 & Act \S3 (CTD) & Act \S4 & Code (a) & Code (b) & Code (c)\\\midrule
Itemised deduction on time & fulfilled & fulfilled & not triggered & fulfilled & violated & not triggered & not triggered\\
Returned late, penalty paid & fulfilled & violated, then repaired & fulfilled & fulfilled & not triggered & not triggered & not triggered\\
Address lost in the post & fulfilled & not triggered & not triggered & fulfilled & not triggered & violated & not triggered\\
Hardship: deposit used as last rent & fulfilled & fulfilled & not triggered & violated & not triggered & not triggered & fulfilled\\
\bottomrule\end{tabular}
}
\caption{Every rule in every history, generated by \texttt{LegalDepositExample.lean}. \S3 is a contrary-to-duty rule: it is triggered only on histories where \S2 was violated.}
\label{tab:legal-status}
\end{table}

\begin{table}[ht]
\centering
\small
% Generated by LegalDepositExample.lean; do not edit.
\begin{tabular}{@{}p{0.2\textwidth}p{0.14\textwidth}p{0.1\textwidth}p{0.12\textwidth}p{0.34\textwidth}@{}}\toprule
History & Under the Act & Under the Code & Landlord knows address from & Normative conflict\\\midrule
Itemised deduction on time & lawful & wrong & day 366 & none\\
Returned late, penalty paid & violated, then repaired & no duty arose & day 366 & none\\
Address lost in the post & lawful & wrong & never & none\\
Hardship: deposit used as last rent & unlawful & right & day 366 & $O^{Act}_L\,G\,\neg\mathit{offset}$ (Act \S4) and $O^{Code}_L\,F\,\mathit{offset}$ (Code c): no future is ideal under both\\
\bottomrule\end{tabular}

\caption{Each history under the Act and under the declared Code, when the landlord first knew the forwarding address, and any normative conflict.}
\label{tab:legal-verdicts}
\end{table}

The hardship history is a genuine dilemma: the Act requires
$G\,\neg\mathit{offset}$ and the Code requires $F\,\mathit{offset}$, so no future is
ideal under both. LeanFM reports the conflict:
\AccessibleEquation{A conflict between codes N1 and N2 for actor a at history h holds when N1 obliges a to ensure psi and N2 obliges a to ensure not psi.}{\[
\mathit{Conflict}(N_1,N_2,a,\psi,h)\iff O^{N_1}_a\psi(h)\wedge O^{N_2}_a\neg\psi(h).
\]}
It does not resolve it. Choosing a priority---law over conscience, or the
reverse---is a stakeholder decision. If one is made, it belongs in the typed
requirements as an explicit ordering, where a reviewer can see it.

\section{What CTL says, and what it does not}

CTL's $A$ and $E$ range over admissible futures, so they cannot by themselves
express ``ought''. They still carry most of the engineering weight once the
normative facts are computed as state predicates, the same technique the
Kerberos discussion used for $\mathit{knows}$ and $\mathit{permitted}$:

\begin{itemize}
\item A deposit registry's own conduct is ordinary safety:
$AG\,\neg(\mathit{receiptIssued}\wedge\neg\mathit{protected})$.
\item Remedies must stay available. Following \emph{possibility means
preparedness}, a violation requires a supported path to a claim:
$AG\,(\mathit{violated}_{\S2}\rightarrow EF\,\mathit{claimFiled})$, with a handling
plan in the requirements.
\item The \emph{late return} history is simultaneously a \emph{counterexample}
for the universal claim that \S2 is always met and a \emph{witness} for
$EF\,\mathit{repaired}$. Checker output should name both roles.
\item Knowledge never exceeds permission, the security invariant, has the same
shape here: a registry must not disclose the tenant's address beyond those
permitted, $AG\,(\mathit{knows}(a,\mathit{address})\rightarrow\mathit{permitted}^{Act}(a,\mathit{address}))$.
\end{itemize}

What CTL does not provide is the per-code quantification over ideal futures,
bounded deadlines measured from a knowledge-triggered start, or the
contrary-to-duty recomputation. The example computes those directly over each
history.

\section{Status}

\path{LegalDepositExample.lean} is an executable book example, not a general
normative checker. It hard-codes the four rules and three code clauses. Turning
it into a LeanFM language feature needs a typed form for a normative code: its
author; each rule's actor, trigger (a state formula that may use
$\mathit{knows}$), deadline, path condition and reparation; and an optional
priority between codes. The checker would return a status per rule per history
with its witness or counterexample trace. Like every LeanFM feature, that form
gets checked requirements and a checked implementation mapping before any
checker code. Regenerate this chapter's tables and diagram with
\texttt{lake env lean --run LegalDepositExample.lean}; the provenance record is
in \path{book/generated/legal-provenance.txt}.
